\section{Shared affine completion rules}
\label{sec:coupled-rules}

The independent-edge optimum is the broadest analytic benchmark.  A
physical construction may require one declared rule across many row
pairs.  We now give the exact finite dual after such a rule is imposed.

\subsection{Global constraint operator}

Let
\[
 \cX_{\cE}=\bigoplus_{e\in\cE}\cX_e
\]
contain one upper-triangle kernel per unordered edge.  A declared affine
rule is encoded by a real-linear map
\[
 \mathscr R:\cX_{\cE}\longrightarrow\cK_{\rm glob}
\]
and a target \(\kappa\in\cK_{\rm glob}\).  The components of
\(\mathscr R\boldsymbol F=\kappa\) include every primitive restriction
\(R_eF_e=k_e\) and may also include:
\begin{itemize}
 \item equality of shared interpolation parameters;
 \item a fixed \(q\)-resolved content-lift relation;
 \item block, side, or scale compatibility;
 \item a linear parametrization shared by one explicitly stacked finite
       coefficient family.
\end{itemize}
Nonlinear or data-adaptive rules are not silently included.

\begin{remark}[Fixed data versus a coefficient class]
\label{rem:fixed-data-coupling}
\Cref{thm:coupled-rule-dual} below fixes one finite target
\(\kappa\) and one positive Schur weight \(c\).  A rule chosen before a
growing or infinite coefficient class introduces an additional
\(\inf_{\rm rule}\sup_{\rm coefficients}\) quantifier.  It is covered
only when the relevant finite coefficient family is explicitly stacked
inside \(\kappa\) and \(\mathscr R\); it is not a consequence of the
single-datum theorem.
\end{remark}

For fixed \(c>0\), define
\begin{align}
 \Theta_c^{\mathscr R}(\kappa)
 =
 \min_{\eta,\boldsymbol F}\quad&\eta
 \label{eq:coupled-primal}\\
 \textup{subject to}\quad&
 \mathscr R\boldsymbol F=\kappa,\notag\\
 &
 \sum_{n\ne m}\frac{c_n}{c_m}
 \|A_{\{m,n\}}F_{\{m,n\}}\|_1
 \le\eta
 \qquad(m\in\cM).\notag
\end{align}

\begin{theorem}[Coupled-rule strong dual]
\label{thm:coupled-rule-dual}
If \(\kappa\in\Ran\mathscr R\), then
\begin{align}
 \Theta_c^{\mathscr R}(\kappa)
 =
 \max\quad&
 [\lambda,\kappa]_{\mathbb R}
 \label{eq:coupled-dual}\\
 \textup{subject to}\quad&
 \alpha_m\ge0,\qquad\sum_m\alpha_m=1,\notag\\
 &
 \bigoplus_{m<n}A_{mn}^*z_{mn}
 =
 \mathscr R^*\lambda,\notag\\
 &
 \|z_{mn}\|_\infty
 \le
 \alpha_m\frac{c_n}{c_m}
 +
 \alpha_n\frac{c_m}{c_n}
 \qquad(m<n).\notag
\end{align}
The primal and dual optima are attained.  If
\(\kappa\notin\Ran\mathscr R\), the primal is infeasible and the dual is
unbounded above.
\end{theorem}

\begin{proof}
Use row multipliers \(\alpha_m\) exactly as in
\cref{thm:fixed-weight-dual}, and use \(\lambda\) for the single global
equality constraint.  The infimum in \(\eta\) imposes
\(\sum_m\alpha_m=1\).  The support-function representation of every
edge norm introduces \(z_{mn}\) with the displayed radius.  The
infimum in the global variable \(\boldsymbol F\) is finite precisely when
the displayed block-adjoint stationarity equation holds.  A feasible
\(\boldsymbol F\) and a strictly larger \(\eta\) give relative Slater
feasibility, so finite-dimensional convex duality gives equality and
dual attainment.

Primal attainment follows separately.  Fix one feasible
\(\boldsymbol F_0\), quotient \(\ker\mathscr R\) by the common zero-cost
subspace
\[
 \left\{
 \boldsymbol G\in\ker\mathscr R:
 A_eG_e=0\text{ for every }e
 \right\},
\]
and choose a linear complement.  On that finite-dimensional complement,
the maximum of the positively weighted row norms controls the variable,
so a minimizing sequence has a convergent representative.  For dual
attainment, \(\alpha\) lies in a compact simplex and every \(z_e\) is
uniformly bounded.  Restricting \(\lambda\) to
\((\ker\mathscr R^*)^\perp\) then makes the adjoint stationarity equation
bound \(\lambda\), and the feasible set is closed.

The infeasible case follows by separating
\(\kappa\) from \(\Ran\mathscr R\), as in
\cref{thm:strong-completion-duality}.
\end{proof}

\begin{corollary}[Exact shared-rule obstruction]
\label{cor:shared-rule-obstruction}
Every feasible tuple
\((\alpha,\lambda,(z_{mn}))\) in
\eqref{eq:coupled-dual} proves that every completion family obeying the
declared shared rule has weighted-Schur cost at least
\[
 [\lambda,\kappa]_{\mathbb R}.
\]
An optimal tuple proves the best possible lower obstruction for that
rule and that fixed \(c\).
\end{corollary}

\begin{remark}[Restriction of scope]
A witness for one \(\mathscr R\) does not obstruct a different completion
class.  Conversely, the independent-edge value can be strictly smaller
than the shared-rule value because its feasible set is larger.  Any
negative conclusion must name the exact rule class that its witness
covers.
\end{remark}

\subsection{Hermitian symmetrization}

Suppose both orientations are initially represented by legal completions.
Define
\[
 \widehat F_{mn}(u,v)
 =
 \frac12
 \left(
 F_{mn}(u,v)+\overline{F_{nm}(v,u)}
 \right)
\]
and reconstruct the reverse by conjugate transpose.

\begin{proposition}[No factor loss from reversal]
\label{prop:Hermitian-symmetrization}
If the two orientation costs are transpose compatible, then
\[
 \cC_D(\widehat F_{mn})
 \le
 \frac12
 \left(
 \cC_D(F_{mn})+\cC_D(F_{nm})
 \right).
\]
The symmetrized pair preserves the common Hermitian primitive datum.
\end{proposition}

\begin{proof}
The primitive restrictions agree by the inherited reversal identity.
The cost is the \(\ell^1\)-norm of a linear analysis map, so convexity and
transpose compatibility give the bound.
\end{proof}

\begin{remark}
The proposition confirms that Hermitian compatibility itself costs no
fixed factor.  Shared \(q\)-provenance or a finite stacked-data rule may
still impose a genuine cross-edge cost because those are additional
constraints, not consequences of reversal.
\end{remark}
